\section{Empirical Predictions and Regional Policy Interpretation}
\label{sec:policy}

The model is intended as a mechanism-based framework rather than as a literal description of any single hub program. Its empirical value lies in clarifying which observable margins correspond to the theoretical objects introduced above, and in identifying predictions that distinguish a hub's visible short-run success from its broader regional welfare effects. For the present paper, the most important implication is that the relevant empirical comparison is often not simply between \emph{hub} and \emph{no hub}, but between a \emph{hub alone} and a \emph{hub embedded in a cross-regional renewal strategy}.

A convenient starting point is the decomposition developed in Section~\ref{sec:short_run}. The hub's private short-run appeal is summarized by
\[
\Pi_H = S_A-F,
\qquad
S_A:=v_LP_A(m_A^1-m_A^0),
\]
while its social value as a stand-alone policy is
\[
\Delta_H=\Pi_H+D_A,
\qquad
D_A:=\beta v_LP_A(m_A^1-m_A^0)\Bigl[1-\delta(m_A^1+m_A^0)\Bigr].
\]
With pipeline policy, social value becomes
\[
\Delta_{H+b}^*=\Pi_H+D_A+R_{AB}(N_A,N_B),
\]
where
\[
R_{AB}(N_A,N_B)=\frac{(\beta\eta v_XQ_{AB})^2}{2\kappa}>0.
\]
This decomposition immediately suggests what should be measured empirically. A good evaluation should not stop at visible activity generated by the hub in region $A$; it should also track whether local interaction becomes repetitive over time and whether policy succeeds in creating cross-regional renewal through outside linkages.

\subsection{Mapping Theory to Observable Margins}

Table \ref{tab:proxy_mapping} summarizes a parsimonious correspondence between the model's theoretical objects and plausible empirical proxies.

\begin{table}[htbp]
\centering
\caption{Illustrative Mapping from Theoretical Objects to Empirical Proxies}
\label{tab:proxy_mapping}
\begin{tabular}{p{0.25\textwidth} p{0.27\textwidth} p{0.38\textwidth}}
\toprule
Theoretical object & Empirical interpretation & Illustrative proxies \\
\midrule
Local matching gain $m_A^1-m_A^0$ 
& Immediate coordination effect of the hub inside region $A$ 
& Interaction volume, first-time meetings, new project starts, participant turnover \\

Dynamic redundancy $\delta$ 
& Speed at which local interaction becomes repetitive or locked in 
& New ties share, repeated collaborator ratio, decline in participant turnover, network modularity or concentration of repeated ties \\

Cross-regional renewal $b$ and $q(b)$ 
& Outside-linkage formation and extra-regional brokerage 
& Outside-region collaborator share, invited outside participants, co-patenting/co-application novelty, external mentor or university linkages \\

Absorptive capacity / anchor presence 
& Capacity to convert contact into productive regional surplus 
& Anchor participation share, lead-firm or university presence, share of boundary-spanning intermediaries \\
\bottomrule
\end{tabular}
\end{table}

The purpose of Table \ref{tab:proxy_mapping} is not to impose a single empirical design. Rather, it highlights that the model's core variables are observable at least in approximate form. In particular, the model suggests that a credible empirical evaluation should include not only gross activity measures, but also indicators of \emph{novelty}, \emph{repetition}, and \emph{outside connectivity}.

\subsection{Three Testable Predictions}

The model yields three testable predictions that are especially relevant for regional policy evaluation.

\paragraph{Prediction 1. Hub creation raises initial interaction volume, but does not necessarily sustain the share of new ties.}

The baseline model implies an immediate local coordination gain,
\[
B_1(1)-B_1(0)=v_LP_A(m_A^1-m_A^0)>0.
\]
Thus, after hub establishment, one should expect an increase in visible local activity in region $A$: more meetings, more matches, more projects, or more observed interaction volume. This is the mechanism behind the private short-run criterion $\Pi_H$.

However, Section~\ref{sec:short_run} also shows that the hub's dynamic correction is
\[
D_A=\beta v_LP_A(m_A^1-m_A^0)\Bigl[1-\delta(m_A^1+m_A^0)\Bigr].
\]
When local redundancy is strong, this term becomes small or negative. In empirical terms, the model therefore predicts that early increases in interaction volume need not translate into persistently high novelty. Instead, one should expect some hubs to exhibit rising \emph{repeated collaborator ratios} or falling \emph{new ties shares} over time, especially where the regional participant pool is narrow and the same actors repeatedly interact.

The empirical implication is straightforward: evaluations based only on attendance, occupancy, or total interactions are incomplete. A more theory-consistent design should track whether first-time collaborations continue to form after the hub's initial launch period. In the model's language, high short-run activity is evidence about $S_A$, but not by itself about $D_A$.

\paragraph{Prediction 2. Regions with a hub plus pipeline policy should exhibit more persistent outside collaboration and higher novelty than otherwise similar hub-only regions.}

Section~\ref{sec:refresh} shows that conditional on a hub, the planner chooses
\[
b^*=\frac{\beta\eta v_XQ_{AB}}{\kappa}>0,
\]
and that the resulting incremental renewal surplus is
\[
R_{AB}(N_A,N_B)=\frac{(\beta\eta v_XQ_{AB})^2}{2\kappa}>0.
\]
The interpretation is that pipeline policy acts on a different margin from the hub itself. The hub raises local coordination inside region $A$, whereas the pipeline policy raises cross-regional connections between region $A$ and the outside pool $B$.

This implies a testable contrast between \emph{hub-only} regions and \emph{hub-plus-pipeline} regions. Relative to otherwise similar hub-only regions, regions that combine the hub with active outside-linkage policy should display a higher \emph{outside-region collaborator share}, greater \emph{co-patenting or co-application novelty}, and a more persistent flow of new collaborative combinations in the medium run. They should also be less likely to exhibit the stagnation pattern described in Prediction 1.

From a policy perspective, this prediction is central. The theory does not merely claim that outside linkages are useful in general. It claims that they are especially valuable when a hub has already increased local interaction intensity and therefore increased the danger of repetition. Empirically, the relevant comparison is therefore not simply before versus after the hub, but whether hub effects differ systematically according to the strength of extra-regional linkage formation.

\paragraph{Prediction 3. The same hub policy generates larger welfare gains in regions with higher absorptive capacity.}

The heterogeneous-participant extension implies that composition matters in addition to scale. A region with the same number of participants may nevertheless generate higher surplus if a larger share of those participants are especially effective at transforming contact into productive collaboration. In practice, this corresponds to the presence of anchor firms, strong universities, capable intermediaries, or other actors that raise the expected value of realized interactions.

The empirical implication is that the effect of hub policy should be stronger in regions with greater \emph{absorptive capacity}. Observable proxies include \emph{anchor participation share}, the presence of universities or lead firms, or the share of participants with boundary-spanning roles. Relative to observationally similar hubs in low-capacity settings, hubs in high-capacity settings should exhibit higher persistence of novel collaborations, stronger project conversion, and larger medium-run gains from outside-linkage policy.

This prediction matters because it sharpens the interpretation of ecosystem thickness. The theory does not say that headcount alone determines efficiency. Rather, it suggests that the effective thickness of a region depends jointly on the number of plausible participants and on the composition of the participant pool. In empirical work, this means that a hub in a smaller region should not automatically be expected to perform poorly if that region contains sufficiently strong anchor actors or unusually capable intermediaries.

\subsection{Implications for Empirical Design}

Taken together, the three predictions imply that a theory-consistent empirical evaluation of hub policy should be dynamic, network-aware, and explicitly regional.

First, evaluation should be \emph{dynamic}. Early activity measures capture the short-run coordination term $S_A$, but do not reveal whether novelty persists. At a minimum, researchers should compare initial interaction gains with the subsequent evolution of new ties share, repeated collaborator ratio, and participant turnover.

Second, evaluation should be \emph{network-aware}. The model is fundamentally about the difference between creating interactions and creating \emph{nonrepetitive} interactions. Measures such as repeated collaborator ratio, outside-region collaborator share, network modularity, and the novelty content of co-patenting or co-application patterns are therefore better aligned with the model than aggregate attendance alone.

Third, evaluation should be \emph{explicitly regional}. The two-region formulation implies that the outside pool is not a residual background feature, but a policy-relevant object. This means that data on external participants, outside mentors, nonlocal co-applicants, or interregional partner flows are not peripheral controls; they are central outcome variables for judging whether a hub is functioning as part of a broader regional strategy.

Overall, the empirical agenda suggested by the model is therefore sharper than a conventional utilization-based evaluation. The critical question is not simply whether the hub generated more activity, but whether it generated activity that remained novel and whether policy design created enough outside renewal to prevent local coordination from hardening into local repetition.
